\documentclass[11pt]{article}
\usepackage[letterpaper,margin=1in]{geometry}
\usepackage{amsmath,amssymb,amsthm,booktabs,microtype,xurl}
\usepackage[hidelinks]{hyperref}
\hypersetup{
  pdftitle={A transversal-capacity obstruction for a regular R(5,5) endpoint},
  pdfauthor={Alec Kriebel},
  pdfsubject={A catalog-conditional endpoint theorem for the Ramsey number R(5,5)},
  pdfkeywords={Ramsey numbers, graph theory, computational combinatorics}
}

\newtheorem{theorem}{Theorem}
\newtheorem{lemma}[theorem]{Lemma}
\newtheorem{corollary}[theorem]{Corollary}
\newcommand{\I}{\mathcal I}
\newcommand{\Q}{\mathcal Q}
\newcommand{\F}{\mathcal F}

\title{A transversal-capacity obstruction for a regular
\texorpdfstring{$R(5,5)$}{R(5,5)} endpoint}
\author{Alec Kriebel\\[2pt]\large with heavy assistance from ChatGPT 5.6 Sol}
\date{Research note prepared 24 July 2026\\Not peer reviewed}

\begin{document}
\maketitle

\begin{abstract}
We give an elementary capacity inequality for hypergraph transversals indexed
by the vertices of a graph.  Applied to the cross-neighborhoods at a vertex of
a hypothetical 18-regular $(5,5;43)$-graph, it excludes the endpoint
\[
 e(G[N(v)])=85.
\]
On byte-pinned published $R(4,5;18)$ and $R(4,5;24)$ catalogs, the inequality
directly excludes 61,939 of 62,382 fixed-side pairs.  The remaining 443 attain
equality.  Uniqueness of a minimizing transversal then forces two adjacent
cross-neighborhoods to coincide, contradicting the required two-column
covering condition.  Conditional on the published catalog-completeness and
extremal-edge statements, every vertex of an 18-regular $(5,5;43)$-graph
therefore lies in at most 84 triangles.  A compact standard-library verifier
reconstructs all finite calculations.  This excludes one endpoint layer; it
does not improve the known bounds on $R(5,5)$.
\end{abstract}

\noindent\textbf{Verification disclaimer.}
This is an unreviewed AI-assisted research note released for expert checking.
The human author is a complete amateur and cannot independently validate the
mathematics.  Exact computation verifies the encoded finite statements; it
does not prove the external catalog-completeness assertions or establish
literature priority.

\section{Introduction}

The diagonal Ramsey number $R(5,5)$ is the least $n$ for which every graph on
$n$ vertices contains either a clique of order five or an independent set of
order five.  Exoo's 42-vertex construction gives the lower bound
$R(5,5)\geq43$~\cite{exoo1989}.  Angeltveit and McKay recently proved
$R(5,5)\leq46$~\cite{angeltveit-mckay-46}.  The exact value remains unknown.

Call a graph with neither a $K_5$ nor an independent five-set a
$(5,5)$-graph.  One possible route to a better upper bound is to assume a
$(5,5)$-graph of order 43, root it at a vertex, and glue its neighborhood to
its anti-neighborhood.  The resulting two sides are smaller $R(4,5)$-graphs.
This strategy has a long history: feasible-cone gluing was central to the
proof that $R(4,5)=25$~\cite{mckay-radziszowski-r45}, and modern
$R(5,5)$ computations combine catalog generation, local gluing, linear
programming, and subgraph-count identities
~\cite{mckay-radziszowski-lp,mckay-radziszowski-counting,
angeltveit-mckay-48,angeltveit-mckay-46}.

The contribution here is narrow.  For each allowed cross-neighborhood size,
we compute the least possible number of independent triples missed by a
cross-neighborhood that meets every independent four-set.  Summing these
minimum-miss penalties across all columns yields a necessary capacity
inequality.  At one regular degree-18 endpoint, the inequality eliminates
almost every fixed-side pair; equality is rigid enough to eliminate the
rest.

\begin{theorem}[Endpoint theorem]\label{thm:endpoint}
Assume the completeness and nonisomorphism statements for the published
edge-85 $R(4,5;18)$ catalog and the complete $R(4,5;24)$ catalog.  No
18-regular $(5,5)$-graph on 43 vertices has a vertex $v$ with
\[
 e(G[N(v)])=85.
\]
Together with the published extremal bound $e(G[N(v)])\leq85$, every vertex
of such a graph lies in at most 84 triangles.
\end{theorem}

The theorem is catalog-conditional and endpoint-local.  It does not exclude
the other regular degree-18 layers, close another global branch, or alter the
known range for $R(5,5)$.

\section{A transversal-capacity inequality}

Let $U$ be a finite set, let $\F,\Q\subseteq2^U$, and let $J$ be a graph with
$\alpha(J)\leq a$.  Assign a set $X_b\subseteq U$ to each $b\in V(J)$ so
that
\begin{enumerate}
\item $X_b$ meets every member of $\F$; and
\item $X_b\cup X_c$ meets every member of $\Q$ whenever $bc\in E(J)$.
\end{enumerate}
For an integer $s$, define
\begin{equation}\label{eq:profile}
 q_s(U;\F,\Q)=
 \min_{\substack{X\subseteq U,\ |X|=s\\
                  X\cap F\ne\varnothing\ \forall F\in\F}}
 \bigl|\{Q\in\Q:X\cap Q=\varnothing\}\bigr|,
\end{equation}
with $q_s=+\infty$ if no size-$s$ transversal of $\F$ exists.

\begin{lemma}[Transversal capacity]\label{lem:capacity}
Every such assignment satisfies
\begin{equation}\label{eq:capacity}
 \sum_{b\in V(J)}q_{|X_b|}(U;\F,\Q)\leq a|\Q|.
\end{equation}
\end{lemma}

\begin{proof}
For $Q\in\Q$, put
\[
 Z_Q=\{b\in V(J):X_b\cap Q=\varnothing\}.
\]
If two vertices of $Z_Q$ were adjacent in $J$, their assigned sets would have
a union disjoint from $Q$, contrary to the second hypothesis.  Thus $Z_Q$ is
independent and $|Z_Q|\leq a$.  Double counting the incidences $(b,Q)$ for
which $X_b\cap Q=\varnothing$ gives
\[
 \sum_bq_{|X_b|}
 \leq\sum_b|\{Q:X_b\cap Q=\varnothing\}|
 =\sum_Q|Z_Q|
 \leq a|\Q|.\qedhere
\]
\end{proof}

\begin{corollary}[Equality rigidity]\label{cor:rigidity}
If equality holds in~\eqref{eq:capacity}, then every $X_b$ attains the
minimum missed-$\Q$ count at its size, and every $Z_Q$ has size $a$.  In
particular, suppose that the size-$s$ minimizer is unique, $q_s>0$, and the
vertices $b$ with $|X_b|=s$ induce an edge in $J$.  Then no assignment
satisfying the two hypotheses of Lemma~\ref{lem:capacity} exists.
\end{corollary}

\begin{proof}
Equality at the endpoints of
\[
 \sum_bq_{|X_b|}
 \leq \sum_Q|Z_Q|
 \leq a|\Q|
\]
forces equality term by term in both inequalities.  Thus each $X_b$
minimizes at its size and every $|Z_Q|=a$.  Under the additional hypotheses,
the columns at the endpoints of an induced edge both equal the unique
minimizer $X^\ast$.  Since $q_s>0$, some $Q\in\Q$ misses $X^\ast$ and hence
also misses the union of those two columns, contradicting the edge
condition.
\end{proof}

The same proof admits fixed nonnegative weights.  If
\[
 q_s^w=
 \min_{\substack{X\subseteq U,\ |X|=s\\
                  X\cap F\ne\varnothing\ \forall F\in\F}}
 \sum_{\substack{Q\in\Q\\X\cap Q=\varnothing}}w_Q,
\]
then
\[
 \sum_bq_{|X_b|}^w\leq a\sum_{Q\in\Q}w_Q.
\]
This weighted statement is an immediate template, not a tested claim about
the remaining $R(5,5)$ layers.

\section{Ramsey specialization}

Let $G$ be a hypothetical 18-regular $(5,5)$-graph on 43 vertices.  Fix
$v\in V(G)$ and put
\[
 A=G[N_G(v)],\qquad
 B=V(G)\setminus(N_G(v)\cup\{v\}),\qquad
 H=\overline{G[B]}.
\]
Then $|A|=18$, $|H|=24$, and both $A$ and $H$ contain neither a $K_4$ nor an
independent five-set.  In particular, $\alpha(H)\leq4$.

For $b\in B=V(H)$ define its cross-neighborhood
\[
 X_b=N_G(b)\cap V(A).
\]
The absence of independent five-sets in $G$ implies two local covering
conditions:
\begin{enumerate}
\item $X_b$ meets every independent four-set of $A$;
\item if $bc\in E(H)$, then $X_b\cup X_c$ meets every independent three-set
of $A$.
\end{enumerate}
Indeed, a missed independent four-set together with $b$, or a missed
independent triple together with the $G$-nonedge $bc$, would be an independent
five-set in $G$.  The second rule is the classical two-column feasible-cone
condition $E_2$ of McKay and Radziszowski
~\cite{mckay-radziszowski-r45}; the first is its immediate one-column
forbidden-set analogue.

Regularity determines every column size:
\[
 18=|X_b|+\deg_{G[B]}(b)=|X_b|+23-d_H(b),
\]
and hence
\begin{equation}\label{eq:column-size}
 |X_b|=d_H(b)-5.
\end{equation}
Let $\I_k(A)$ denote the independent $k$-sets of $A$, put
$i_3(A)=|\I_3(A)|$, and specialize~\eqref{eq:profile} to
\begin{equation}\label{eq:ramsey-profile}
 q_s(A)=
 \min_{\substack{X\subseteq V(A),\ |X|=s\\
 X\cap I\ne\varnothing\ \forall I\in\I_4(A)}}
 |\{Q\in\I_3(A):X\cap Q=\varnothing\}|.
\end{equation}
Lemma~\ref{lem:capacity}, with $J=H$ and $a=4$, gives the necessary condition
\begin{equation}\label{eq:ramsey-capacity}
 \boxed{\displaystyle
 \sum_{b\in V(H)}q_{d_H(b)-5}(A)\leq4i_3(A).}
\end{equation}

The two fixed-side edge counts are linked exactly.  Summing degrees over
$A$ shows that the number of $A$--$B$ edges is
\[
 18\cdot18-18-2e(A)=306-2e(A).
\]
On the $B$ side,~\eqref{eq:column-size} gives
\[
 \sum_{b\in B}|X_b|=\sum_{b\in V(H)}(d_H(b)-5)=2e(H)-120.
\]
Equating these quantities yields
\begin{equation}\label{eq:edge-balance}
 e(H)=213-e(A).
\end{equation}
Thus $e(A)=85$ is exactly the endpoint $e(H)=128$.

\section{The finite endpoint calculation}

The external catalog inputs are the publisher-supplied graph6 files on
McKay's Ramsey data page~\cite{mckay-data}.  Angeltveit and McKay's 2026
paper records $E(4,5;18)=85$, the 74 edge-85 order-18 graphs, and the 843
edge-128 order-24 graphs~\cite{angeltveit-mckay-46}.  The relevant data files
contain those 74 endpoint records and the complete 352,366-record
$R(4,5;24)$ catalog.  Completeness of the order-24 catalog is Theorem 2.1 of
Angeltveit--McKay~\cite{angeltveit-mckay-48}.

The exact input bytes used here are
\begin{align*}
\operatorname{SHA256}(\texttt{r4518.85.g6})
 &=\texttt{46abaee2572d06bba1e594554809d784be}\\[-2pt]
 &\quad\texttt{60f8f60b9b0d3345b8bf3dd800810a},\\
\operatorname{SHA256}(\texttt{r45\_24.g6})
 &=\texttt{83ca4028f206b2fa4315ef219b8c2c57}\\[-2pt]
 &\quad\texttt{c7835209673dd8183d8fb4353bd4fdd0}.
\end{align*}

For every order-18 record $A$, the verifier enumerates all independent triples
and four-sets, then every exact-size subset needed by the degree sequences of
the selected $H$ records.  It retains precisely those subsets that meet every
independent four-set and computes their missed-triple counts.  This determines
the exact profile values~\eqref{eq:ramsey-profile}.  Each of the
$74\cdot843=62,382$ fixed-side pairs is then classified by
\eqref{eq:ramsey-capacity}.
The profiles, degree data, and terminal equality conditions are invariant
under relabeling.  Thus every choice of side isomorphisms is covered; no
cross-side vertex matching is an omitted parameter.

\begin{center}
\begin{tabular}{lr}
\toprule
Classification & Pair count\\
\midrule
Strict violation of the capacity inequality & 61,939\\
Equality cases requiring terminal analysis & 443\\
Pairs admitting a completion after terminal analysis & 0\\
\midrule
Total fixed-side pairs & 62,382\\
\bottomrule
\end{tabular}
\end{center}

The 61,939 strict violations cannot admit any cross-edge completion.  The next
section eliminates every equality pair without enumerating cross-edge
matrices.

\begin{theorem}[Finite pinned-catalog theorem]\label{thm:finite}
Let $\mathcal A_{85}$ be the 74 graphs encoded in the pinned edge-85
$R(4,5;18)$ endpoint file, and let $\mathcal H_{128}$ be the 843 edge-128
graphs selected from the pinned $R(4,5;24)$ file.  For no pair
\[
 (A,H)\in\mathcal A_{85}\times\mathcal H_{128}
\]
does there exist a family $(X_b)_{b\in V(H)}$ satisfying
the prescribed sizes~\eqref{eq:column-size}, the one-column
independent-four-set condition, and the two-column independent-triple
condition.
\end{theorem}

\section{Equality rigidity}

All 443 equality pairs use the same order-18 graph $A$, at zero-based catalog
index 50, and every associated $H$ has degree sequence
\[
 10^8\,11^{16}.
\]
For this exceptional $A$, exact enumeration gives
\[
 i_3(A)=74,\qquad q_5(A)=17,\qquad q_6(A)=10.
\]
Moreover, the size-six minimizer in~\eqref{eq:ramsey-profile} is unique; call
it $X^\ast$.  The capacity inequality is tight:
\begin{equation}\label{eq:tight}
 8q_5(A)+16q_6(A)
 =8\cdot17+16\cdot10
 =296
 =4i_3(A).
\end{equation}

Suppose a cross-neighborhood assignment existed.  The proof of
Lemma~\ref{lem:capacity} gives a chain
\[
 \sum_bq_{|X_b|}(A)
 \leq\sum_{Q\in\I_3(A)}|Z_Q|
 \leq4i_3(A).
\]
Equation~\eqref{eq:tight} forces equality throughout.  By
Corollary~\ref{cor:rigidity}, every column attains its individual profile
minimum.  In particular, every one of the sixteen degree-11 vertices of $H$
has a size-six cross-neighborhood equal to the unique minimizer $X^\ast$.

Let $D$ be those sixteen vertices.  Each $b\in D$ has at most eight neighbors
outside $D$, so it has at least three neighbors in $D$.  Thus $H[D]$ contains
an edge $bc$.  But
\[
 X_b\cup X_c=X^\ast,
\]
and $X^\ast$ misses ten independent triples of $A$.  This contradicts the
two-column condition on the edge $bc$.  All 443 equality pairs are therefore
impossible, proving Theorem~\ref{thm:finite}.  The published catalog
coverage and extremal-edge statements then give Theorem~\ref{thm:endpoint}.

\section{Verification and reproducibility}

The computation is deterministic, solver-free, and uses only Python's
standard library.  The public bundle contains:
\begin{itemize}
\item the two byte-pinned source catalogs and their license attribution;
\item the certificate producer and a separately implemented checker that
imports no producer code;
\item the complete 62,382-line classification stream;
\item 17 focused semantic tests; and
\item a one-command launcher that verifies file hashes, copies the release to
a temporary directory, regenerates the frozen artifacts byte-for-byte, runs
the checker, and runs the tests.
\end{itemize}

The corrected v1.0.1 normalized release archive has SHA-256
\begin{center}
\texttt{de541d6c7ed8be496784397ea0ee3f1}\\[-2pt]
\texttt{b12c2b93cdbc42ba908160095c1d79cc4}.
\end{center}
A fresh copy downloaded from the public release returned
\texttt{"valid": true}, reproduced every frozen artifact, and passed all 17
tests in under 70 seconds with peak resident memory below 140 MB.  The release,
source branch, and verification instructions are archived at
\begin{center}
\url{https://github.com/AlecKriebel/Math/releases/tag/ramsey55-endpoint-capacity-v1.0.1}.
\end{center}

The verifier checks the exact supplied bytes, graph validity, Ramsey
properties, edge filters, profile values, pair classifications, and equality
witnesses.  It does not independently regenerate the complete
$R(4,5)$ catalogs.  Catalog completeness and pairwise nonisomorphism remain
external hypotheses.

\section{Relation to prior work and scope}

The two-column covering rule is the $E_2$ feasible-cone condition used by
McKay and Radziszowski in their computation of
$R(4,5)$~\cite{mckay-radziszowski-r45}; the one-column rule is its immediate
forbidden-set analogue.  The summation in Lemma~\ref{lem:capacity} is an
elementary incidence double count.  Global
neighborhood identities and linear-programming aggregation also have a long
history in Ramsey computation
~\cite{mckay-radziszowski-lp,mckay-radziszowski-counting}.
Gauthier and Brown have since given a kernel-checked HOL4 proof of
$R(4,5)=25$ using verified gluing lemmas~\cite{gauthier-brown}.

The defensible candidate contribution is more specific: we formulate the
exact
minimum-miss profile $q_s(A)$, its aggregate use across graph-indexed columns,
and the equality-rigidity argument that closes the $(85,128)$ endpoint.  A
focused primary-source search found no earlier statement of this profile
inequality or endpoint theorem.  That search cannot establish worldwide
priority; unindexed code, unpublished filters, or differently phrased
antecedents may exist.

No claim is made that the abstract double count is itself new, that this is
the first use of hypergraph transversals in Ramsey computation, or that the
catalog-conditional endpoint theorem changes $R(5,5)$.  The remaining regular
degree-18 layers and five other global branches remain open.  A meaningful
continuation would require a weighted or correlated multi-column capacity
inequality with whole-layer reach, a global counting theorem forcing a vertex
into this endpoint, or a comparably strong new construction.

\appendix
\section{Exceptional endpoint data}

For independent reproduction of the hand-checkable terminal step, the
exceptional order-18 record is catalog line 51 (zero-based index 50):
\begin{center}
\verb|QznZZYVk{mHZuJuD@\XfQOs}_lo|
\end{center}
Its unique size-six minimizer is
\[
 X^\ast=\{2,3,4,9,11,12\}
\]
in zero-based vertex numbering.  The ten independent triples missed by
$X^\ast$ are
\[
\begin{gathered}
\{0,6,7\},\ \{0,14,15\},\ \{1,8,10\},\ \{1,14,15\},\
\{5,14,15\},\\
\{5,16,17\},\ \{6,7,13\},\ \{6,7,17\},\
\{8,10,13\},\ \{8,10,16\}.
\end{gathered}
\]
These labeled data are not additional hypotheses: the separately implemented
checker reconstructs them from the pinned graph6 record.  The catalog files
are attributed under the publisher's stated CC BY 4.0 data license.

\section*{AI-assistance and authorship disclosure}

The exploratory mathematics, proof drafting, programs, independent-code
audits, literature review, and publication materials were developed with
heavy assistance from ChatGPT 5.6 Sol under Alec Kriebel's direction.  Alec
Kriebel is the named human author.  No external expert has reviewed this note,
and passing the supplied verifier is not peer review.

\begin{thebibliography}{99}

\bibitem{exoo1989}
G.~Exoo,
\newblock A lower bound for $R(5,5)$,
\newblock \emph{Journal of Graph Theory} 13 (1989), 97--98.
\newblock \url{https://doi.org/10.1002/jgt.3190130113}

\bibitem{mckay-radziszowski-lp}
B.~D. McKay and S.~P. Radziszowski,
\newblock Linear programming in some Ramsey problems,
\newblock \emph{Journal of Combinatorial Theory, Series B} 61 (1994),
125--132.
\newblock \url{https://doi.org/10.1006/jctb.1994.1038}

\bibitem{mckay-radziszowski-r45}
B.~D. McKay and S.~P. Radziszowski,
\newblock $R(4,5)=25$,
\newblock \emph{Journal of Graph Theory} 19 (1995), 309--322.
\newblock \url{https://doi.org/10.1002/jgt.3190190304}

\bibitem{mckay-radziszowski-counting}
B.~D. McKay and S.~P. Radziszowski,
\newblock Subgraph counting identities and Ramsey numbers,
\newblock \emph{Journal of Combinatorial Theory, Series B} 69 (1997),
193--209.
\newblock \url{https://doi.org/10.1006/jctb.1996.1741}

\bibitem{angeltveit-mckay-48}
V.~Angeltveit and B.~D. McKay,
\newblock $R(5,5)\leq48$,
\newblock \emph{Journal of Graph Theory} 89 (2018), 5--13.
\newblock \url{https://doi.org/10.1002/jgt.22235}

\bibitem{gauthier-brown}
T.~Gauthier and C.~E. Brown,
\newblock A formal proof of $R(4,5)=25$,
\newblock in \emph{15th International Conference on Interactive Theorem
Proving}, LIPIcs 309 (2024), 16:1--16:18.
\newblock \url{https://doi.org/10.4230/LIPIcs.ITP.2024.16}

\bibitem{angeltveit-mckay-46}
V.~Angeltveit and B.~D. McKay,
\newblock $R(5,5)\leq46$,
\newblock \emph{Journal of Graph Theory} 112 (2026), 198--208.
\newblock \url{https://doi.org/10.1002/jgt.70029}

\bibitem{mckay-data}
B.~D. McKay,
\newblock Ramsey graph data,
\newblock accessed 24 July 2026.
\newblock \url{https://users.cecs.anu.edu.au/~bdm/data/ramsey.html}

\end{thebibliography}

\end{document}
